\documentclass[a4paper,11pt]{article}
\usepackage[T2A]{fontenc}
\usepackage[utf8]{inputenc}
\usepackage[russian]{babel}
\usepackage{amsmath,amssymb,mathtools}
\usepackage{geometry}
\geometry{left=21mm,right=21mm,top=22mm,bottom=23mm}
\usepackage{microtype}
\usepackage{icomma}
\usepackage{tikz}
\usetikzlibrary{arrows.meta,positioning,shapes.geometric,calc}
\usepackage{tkz-euclide}
\usepackage{tabularx,booktabs,longtable,array}
\usepackage{enumitem}
\usepackage{fancyhdr}
\usepackage{setspace}
\usepackage{hyperref}
\usepackage{xcolor}
\definecolor{accent}{HTML}{087D82}
\definecolor{darkaccent}{HTML}{005B61}
\definecolor{textgray}{HTML}{353535}
\definecolor{muted}{HTML}{686868}
\hypersetup{colorlinks=true,linkcolor=darkaccent,urlcolor=darkaccent,pdfauthor={Лёвин Артём Александрович},pdftitle={Чек-лист: четырёхугольники и квадратные корни — Екатерина Скелли}}
\setstretch{1.19}
\setlength{\parindent}{0pt}
\setlength{\parskip}{5pt}
\setlist[itemize]{leftmargin=1.7em,itemsep=2pt,topsep=3pt}
\setlist[enumerate]{leftmargin=1.85em,itemsep=5pt,topsep=3pt}
\pagestyle{fancy}
\fancyhf{}
\fancyhead[L]{\small\color{muted}Четырёхугольники и квадратные корни}
\fancyhead[R]{\small\color{muted}Екатерина Скелли}
\fancyfoot[C]{\small\thepage}
\renewcommand{\headrulewidth}{0.45pt}
\renewcommand{\headrule}{\hbox to\headwidth{\color{accent}\leaders\hrule height\headrulewidth\hfill}}
\renewcommand{\footrulewidth}{0pt}
\setlength{\headheight}{22pt}
\newcommand{\sect}[1]{\vspace{7pt}\section*{\color{darkaccent}#1}\vspace{-5pt}\noindent{\color{accent}\rule{\linewidth}{0.45pt}}\vspace{3pt}}
\newcommand{\sub}[1]{\vspace{5pt}{\color{darkaccent}\large #1}\par\vspace{0pt}}
\newcommand{\key}[1]{\textcolor{darkaccent}{#1}}
\newcommand{\ruleline}{\par\vspace{2pt}\noindent{\color{accent}\rule{\linewidth}{0.45pt}}\par}
\begin{document}
\thispagestyle{empty}
\begin{titlepage}
  \centering
  \vspace*{23mm}
  {\Large\color{darkaccent}\textls[65]{МАТЕМАТИКА}\par}
  \vspace{25mm}
  {\Huge Чек-лист\par}
  \vspace{12mm}
  {\LARGE\color{accent}Четырёхугольники\\[3mm]и квадратные корни\par}
  \vspace{21mm}
  {\large Повторение и самостоятельная практика\\[2mm]8 класс\par}
  \vfill
  {\Large Екатерина Скелли\par}
  \vspace{13mm}
  {\large 17 ноября 2026 г.\par}
  \vspace{18mm}
  {\normalsize Лёвин Артём Александрович\\[1mm]эксклюзивно для Екатерины Скелли\par}
  \vspace{18mm}
  \begin{tikzpicture}[x=1cm,y=1cm]
    \draw[accent,line width=.9pt] (-4.1,0) .. controls (-2.9,.68) and (-1.9,-.48) .. (0,0)
      .. controls (1.7,.62) and (2.5,-.57) .. (4.1,0);
    \fill[darkaccent] (0,0) circle (1.2pt);
  \end{tikzpicture}
  \vspace{9mm}
\end{titlepage}

\sect{1. Прямоугольник: диагонали и углы}
\textbf{Проверьте, что Вы знаете:}
\begin{itemize}
\item противоположные стороны прямоугольника равны, каждый угол равен $90^\circ$;
\item диагонали прямоугольника \textbf{равны} и \textbf{делятся точкой пересечения пополам};
\item если $AC$ и $BD$ пересекаются в $O$, то $AO=BO=CO=DO$;
\item в равнобедренном треугольнике углы при основании равны;
\item сумма углов треугольника $180^\circ$, смежные углы также дают $180^\circ$.
\end{itemize}

\sub{Пример 1. Угол между диагоналями}
В прямоугольнике $ABCD$ диагонали пересекаются в точке $O$; $\angle ACD=38^\circ$. Найти $\angle BOC$.
\begin{center}
\begin{tikzpicture}[scale=.72]
  \tkzDefPoints{0/2.5/A,5/2.5/B,5/0/C,0/0/D,2.5/1.25/O}
  \tkzDrawPolygon[thick](A,B,C,D)
  \tkzDrawSegments[accent,thick](A,C B,D)
  \tkzDrawPoints(O)
  \tkzLabelPoints[above left](A)
  \tkzLabelPoints[above right](B)
  \tkzLabelPoints[below right](C)
  \tkzLabelPoints[below left](D)
  \tkzLabelPoints[above](O)
  \node[font=\small] at (4.35,.23) {$38^\circ$};
\end{tikzpicture}
\end{center}
Так как $DO=CO$, треугольник $DOC$ равнобедренный. Поэтому
\[
\angle ODC=\angle OCD=38^\circ,\qquad
\angle DOC=180^\circ-38^\circ-38^\circ=104^\circ.
\]
Углы $DOC$ и $BOC$ смежные. Следовательно,
\[
\boxed{\angle BOC=180^\circ-104^\circ=76^\circ.}
\]
\key{Главная идея:} равные половины диагоналей позволяют получить равнобедренный треугольник.

\sub{Дополнительный приём с точки $K$}
Если на стороне $BC$ отмечена точка $K$ и $BK=BO$, а $\angle BOK=67^\circ$, то в равнобедренном треугольнике $BOK$:
\[
\angle OBK=180^\circ-2\cdot67^\circ=46^\circ.
\]
Диагонали прямоугольника образуют со стороной $BC$ одинаковые острые углы; следовательно,
\[\boxed{\angle BOC=180^\circ-2\cdot46^\circ=88^\circ.}\]

\clearpage
\sect{2. Биссектриса прямого угла}
Биссектриса делит угол $90^\circ$ на два угла по $45^\circ$. В прямоугольнике биссектриса угла $A$, достигающая стороны $BC$ в точке $M$, образует прямоугольный равнобедренный треугольник $ABM$.

\sub{Пример 2. Периметр прямоугольника}
Дано: $BM=7$ см и $MC=4$ см. Найти периметр $ABCD$.
\begin{center}
\begin{tikzpicture}[scale=.67]
 \tkzDefPoints{0/0/A,0/3/B,5/3/C,5/0/D,3/3/M}
 \tkzDrawPolygon[thick](A,B,C,D)
 \tkzDrawSegment[accent,thick](A,M)
 \tkzDrawPoints(M)
 \tkzLabelPoints[below left](A)
 \tkzLabelPoints[above left](B)
 \tkzLabelPoints[above right](C)
 \tkzLabelPoints[below right](D)
 \tkzLabelPoints[above](M)
 \node[above,font=\small] at (1.5,3.1) {7};
 \node[above,font=\small] at (4,3.1) {4};
\end{tikzpicture}
\end{center}
В треугольнике $ABM$ углы при $A$ и $M$ равны $45^\circ$, значит $AB=BM=7$ см. Кроме того, $BC=7+4=11$ см. Тогда
\[\boxed{P=2(AB+BC)=2(7+11)=36\text{ см}.}\]
\key{Не перепутайте:} диагональ прямоугольника и биссектриса угла~--- разные отрезки.

\clearpage
\sect{3. Ромб: диагонали и высота}
\textbf{Набор свойств:}
\begin{itemize}
\item у ромба четыре равные стороны; противоположные углы равны, соседние в сумме дают $180^\circ$;
\item диагонали ромба перпендикулярны и делят его углы пополам;
\item в прямоугольном треугольнике катет напротив угла $30^\circ$ равен половине гипотенузы;
\item верно и обратное: если катет равен половине гипотенузы, против него лежит угол $30^\circ$.
\end{itemize}

\sub{Пример 3. Меньшая диагональ}
Сторона ромба равна $11$ см, один угол равен $120^\circ$. Найти меньшую диагональ и углы, которые она образует со сторонами ромба.
\begin{center}
\begin{tikzpicture}[scale=.75]
 \tkzDefPoints{-3/0/A,0/1.73/B,3/0/C,0/-1.73/D}
 \tkzDrawPolygon[thick](A,B,C,D)
 \tkzDrawSegments[accent,thick](B,D A,C)
 \tkzLabelPoints[left](A)
 \tkzLabelPoints[above](B)
 \tkzLabelPoints[right](C)
 \tkzLabelPoints[below](D)
 \node[font=\small] at (.5,1.0) {$60^\circ$};
 \node[font=\small] at (-.52,1.0) {$60^\circ$};
\end{tikzpicture}
\end{center}
Если $\angle B=120^\circ$, то $\angle A=60^\circ$. В треугольнике $ABD$ стороны $AB=AD=11$ см, а угол между ними $60^\circ$. Значит, треугольник $ABD$ равносторонний и
\[\boxed{BD=11\text{ см}.}\]
Диагональ $BD$ делит углы $B$ и $D$ пополам: $120^\circ:2=60^\circ$.
\[\boxed{\text{Углы меньшей диагонали со сторонами: }60^\circ.}\]
\textit{Второй путь:} половина меньшей диагонали равна $5{,}5$ см как катет против угла $30^\circ$ в одном из прямоугольных треугольников.

\sub{Пример 4. Высота ромба}
В ромбе $ABCD$ угол $A$ острый, $BH\perp AD$, $AB=2\,BH$. Найти углы ромба.
Треугольник $ABH$ прямоугольный, гипотенуза $AB$ вдвое больше катета $BH$. Значит,
\[
\angle BAH=30^\circ,\qquad
\boxed{\angle A=\angle C=30^\circ,\quad \angle B=\angle D=150^\circ.}
\]
\key{Проверка:} $30^\circ+150^\circ=180^\circ$.

\clearpage
\sect{4. Квадрат и равносторонний треугольник}
В квадрате $ABCD$ точка $K$ расположена внутри так, что треугольник $ABK$ равносторонний. Найти $\angle DKC$.
\begin{center}
\begin{tikzpicture}[scale=.64]
 \tkzDefPoints{0/0/A,4/0/B,4/4/C,0/4/D,2/3.464/K}
 \tkzDrawPolygon[thick](A,B,C,D)
 \tkzDrawSegments[accent,thick](A,K B,K D,K C,K)
 \tkzDrawPoints(K)
 \tkzLabelPoints[below left](A)
 \tkzLabelPoints[below right](B)
 \tkzLabelPoints[above right](C)
 \tkzLabelPoints[above left](D)
 \tkzLabelPoints[below](K)
\end{tikzpicture}
\end{center}
Поскольку $AD=AB=AK$, треугольник $ADK$ равнобедренный. Угол $DAK=90^\circ-60^\circ=30^\circ$, поэтому $\angle DKA=(180^\circ-30^\circ)/2=75^\circ$. Аналогично $\angle BKC=75^\circ$. Вокруг точки $K$ сумма углов равна $360^\circ$:
\[
\boxed{\angle DKC=360^\circ-75^\circ-60^\circ-75^\circ=150^\circ.}
\]

\vspace{5mm}
\sect{5. Что означает арифметический квадратный корень}
\textbf{Определение.} Для $a\ge0$ число $\sqrt a$~--- \textbf{неотрицательное} число, квадрат которого равен $a$.
\[\sqrt{a}\ge0,\qquad (\sqrt a)^2=a\quad(a\ge0).\]
\begin{center}
\begin{tabular}{@{}llllll@{}}
\toprule
$\sqrt0=0$ & $\sqrt1=1$ & $\sqrt4=2$ & $\sqrt9=3$ & $\sqrt{25}=5$ & $\sqrt{100}=10$\\
\bottomrule
\end{tabular}
\end{center}
\textbf{Важно:} $\sqrt{-4}$ в действительных числах не определён. Кроме того, $\sqrt{a^2}=|a|$; например, $\sqrt{(-5)^2}=5$.

\sect{6. Произведение и частное корней}
Для $a\ge0$ и $b\ge0$:
\[\boxed{\sqrt a\cdot\sqrt b=\sqrt{ab}.}\]
Для $a\ge0$, $b>0$:
\[\boxed{\frac{\sqrt a}{\sqrt b}=\sqrt{\frac ab}.}\]
Оба равенства применимы и в обратном направлении. 
\[\sqrt{81\cdot100}=\sqrt{81}\,\sqrt{100}=9\cdot10=90.\]
\textbf{Нельзя переносить знак корня через сложение:}
\[\sqrt a+\sqrt b\ne\sqrt{a+b}\quad\text{в общем случае}.\]
Например, $\sqrt9+\sqrt{16}=3+4=7$, тогда как $\sqrt{9+16}=5$.

\sub{Пример 5. Десятичные дроби под корнем}
\[
\sqrt{16{,}9}\cdot\sqrt{0{,}4}
=\sqrt{16{,}9\cdot0{,}4}
=\sqrt{6{,}76}
=\sqrt{\frac{676}{100}}
=\frac{26}{10}
=\boxed{2{,}6}.
\]
При работе с десятичными дробями удобно сначала получить обыкновенную дробь с полным квадратом в числителе и знаменателе.

\sub{Пример 6. Вычисляем корни по отдельности}
\[
\begin{aligned}
0{,}2\sqrt{900}+1{,}8\sqrt{\frac19}
&=0{,}2\cdot30+1{,}8\cdot\frac13\\
&=6+0{,}6=\boxed{6{,}6}.
\end{aligned}
\]
Здесь полезны равенства $\sqrt{900}=\sqrt9\sqrt{100}=30$ и $\sqrt{1/9}=1/3$.

\clearpage
\sect{7. Разложение подкоренного выражения на множители}
\textbf{Маршрут:} объедините умножаемые корни $\to$ разложите числа на простые множители $\to$ соберите одинаковые множители в пары $\to$ вынесите их из-под корня.

\sub{Пример 7. Несколько корней сразу}
Вычислите $\sqrt6\cdot\sqrt{21}\cdot\sqrt{42}\cdot\sqrt{147}$.
\[
\begin{aligned}
&\sqrt{6\cdot21\cdot42\cdot147}\\
&=\sqrt{(2\cdot3)(3\cdot7)(2\cdot3\cdot7)(3\cdot7^2)}\\
&=\sqrt{2^2\cdot3^4\cdot7^4}
=2\cdot3^2\cdot7^2\\
&=2\cdot9\cdot49=\boxed{882}.
\end{aligned}
\]
\key{Смысл «пары»:} из $\sqrt{p^2}$ выходит один множитель $p$ при $p\ge0$. Например, $\sqrt{2^2\cdot3^4}=2\cdot3^2=18$.

\sect{8. Ошибки, которые важно исключить}
\begin{enumerate}[label=\textcolor{accent}{\arabic*.}]
\item Сложение под одним корнем: $\sqrt a+\sqrt b\ne\sqrt{a+b}$.
\item Подмена арифметического корня двумя числами: $\sqrt{25}=5$; уравнение $x^2=25$ имеет \textbf{два} решения, $x=\pm5$.
\item Неверная обработка отрицательного числа: $\sqrt{-9}$ в $\mathbb R$ не существует.
\item Потеря условия $b>0$ в частном $\sqrt a/\sqrt b$.
\item Необоснованная замена $\sqrt{a^2}$ на $a$ при неизвестном знаке $a$.
\item Перепутанные диагональ, биссектриса и высота; неотмеченный прямой угол.
\item Пропуск единиц измерения и знака градусов в геометрическом ответе.
\end{enumerate}

\clearpage
\thispagestyle{fancy}
\begin{center}
{\Large\color{darkaccent}Алгоритм: как работать с квадратными корнями}\par
\vspace{6mm}
\end{center}
\begin{center}
\begin{tikzpicture}[
 node distance=12mm and 13mm,>=Stealth,
 start/.style={ellipse,draw=accent,very thick,align=center,text width=49mm,minimum height=10mm,inner sep=3pt},
 action/.style={rectangle,rounded corners=2pt,draw=accent,align=center,text width=61mm,minimum height=14mm,inner sep=4pt},
 choice/.style={diamond,aspect=2.1,draw=darkaccent,very thick,align=center,text width=43mm,inner sep=2pt},
 flow/.style={-{Stealth[length=2mm]},thick,draw=darkaccent},
 lab/.style={font=\small,text=darkaccent,fill=white,inner sep=1.5pt}]
 \node[start] (s) {Выражение с квадратными корнями};
 \node[action,below=of s] (domain) {Проверьте: подкоренные выражения неотрицательны; знаменатели положительны там, где делят на корень};
 \node[choice,below=of domain] (q) {Корни умножаются или делятся?};
 \node[action,below=20mm of q] (yes) {Объедините под одним корнем; разложите на множители; выделите полные квадраты};
 \node[action,right=16mm of yes,text width=42mm] (no) {При сумме или разности упростите каждый корень отдельно};
 \node[action,below=of yes] (calc) {Вычислите выражение и проверьте результат};
 \node[start,below=of calc] (end) {Ответ};
 \draw[flow] (s)--(domain);
 \draw[flow] (domain)--(q);
 \draw[flow] (q)--node[lab,left]{да}(yes);
 \draw[flow] (q.east) -| node[lab,near start,above]{нет} (no.north);
 \draw[flow] (yes)--(calc);
 \draw[flow] (no.south) |- (calc.east);
 \draw[flow] (calc)--(end);
\end{tikzpicture}
\end{center}
\vspace{5mm}
\begin{center}\textit{Выбирайте ветку по действию между корнями.\\Перед вычислением всегда проверяйте допустимость выражения.}\end{center}

\clearpage
\sect{Тренировочный блок: путь А}
\textbf{10 заданий.} Выполняйте по порядку; в геометрии обозначайте существенные элементы чертежа и обязательно записывайте единицы измерения или знак $^\circ$.
\sub{Четырёхугольники}
\begin{enumerate}[label=\textbf{\arabic*.}]
\item В прямоугольнике $ABCD$ диагонали пересекаются в точке $O$. Известно, что $\angle ACD=32^\circ$. Найдите $\angle BOC$.
\item Биссектриса угла $A$ прямоугольника $ABCD$ пересекает $BC$ в точке $M$. $BM=6$ см, $MC=9$ см. Найдите периметр прямоугольника.
\item Сторона ромба равна $14$ см, один из углов $120^\circ$. Найдите меньшую диагональ ромба и углы между нею и сторонами.
\item В ромбе $ABCD$ угол $A$ острый, высота $BH$ проведена к стороне $AD$ и $AB=2\,BH$. Найдите все углы ромба.
\item В квадрате $ABCD$ точка $K$ находится внутри, треугольник $ABK$ равносторонний. Найдите угол $DKA$.
\end{enumerate}
\sub{Квадратные корни}
\begin{enumerate}[label=\textbf{\arabic*.},start=6]
\item Вычислите: $\sqrt{12{,}25}\cdot\sqrt{0{,}16}$.
\item Вычислите: $0{,}4\sqrt{625}+1{,}5\sqrt{\frac49}$.
\item Вычислите: $\sqrt{14400}$.
\item Вычислите: $\displaystyle\frac{\sqrt{12}\cdot\sqrt{75}}{\sqrt9}$.
\item Вычислите, используя разложение на простые множители: $\sqrt8\cdot\sqrt{14}\cdot\sqrt{21}\cdot\sqrt3$.
\end{enumerate}
\ruleline
\textbf{Самопроверка перед сдачей:} в решении есть обоснование свойства, арифметика проверена, в ответе указаны градусы/сантиметры, применимость свойств корней соблюдена.

\clearpage
\sect{Ответы к тренировочному блоку}
\renewcommand{\arraystretch}{1.48}
\begin{center}
\begin{tabularx}{\linewidth}{@{}>{\bfseries}p{12mm}X@{}}
\toprule
№ & Ответ\\
\midrule
1 & $64^\circ$\\
2 & $42$ см\\
3 & Меньшая диагональ $14$ см; углы $60^\circ$\\
4 & $30^\circ$, $150^\circ$, $30^\circ$, $150^\circ$\\
5 & $75^\circ$\\
6 & $1{,}4$\\
7 & $11$\\
8 & $120$\\
9 & $10$\\
10 & $84$\\
\bottomrule
\end{tabularx}
\end{center}
\vspace{7mm}
\sect{Итоговая памятка}
\begin{tabularx}{\linewidth}{@{}p{.36\linewidth}X@{}}
\toprule
\textbf{Ситуация} & \textbf{Первое полезное действие}\\
\midrule
Диагонали прямоугольника & Отметьте четыре равные половины диагоналей\\
Биссектриса угла $90^\circ$ & Отметьте два угла по $45^\circ$\\
Диагональ ромба & Проверьте деление угла пополам и перпендикулярность диагоналей\\
Катет вдвое короче гипотенузы & Отметьте угол $30^\circ$ напротив этого катета\\
Произведение или частное корней & Объедините корни при допустимых значениях\\
Много множителей под корнем & Разложите на простые множители и объедините в пары\\
Сумма корней & Упростите каждый корень отдельно\\
\bottomrule
\end{tabularx}
\vfill
\begin{center}\small\color{muted}Лёвин Артём Александрович --- эксклюзивно для Екатерины Скелли\end{center}
\end{document}
